\documentclass[11pt]{article}
\usepackage[UTF8]{ctex}
\usepackage{amsmath,amssymb,amsthm}
\usepackage{etoolbox}

% 1. 自定义定理样式实现标题后强制换行
\newtheoremstyle{break}% 样式名称
  {\topsep}{\topsep}% 环境上下间距
  {\normalfont}% 正文字体
  {}% 缩进量
  {\bfseries}{.}% 头部字体和标点
  {\newline}% 头部后的间距：设置为 \newline 即可实现完美的换行效果
  {}% 头部时间/附加参数规范

% 应用新样式并声明常用的定理类环境
\theoremstyle{break}
\newtheorem{theorem}{定理}[section]
\newtheorem{lemma}[theorem]{引理}
\newtheorem{corollary}[theorem]{推论}
\newtheorem{definition}[theorem]{定义}

% 2. 利用 etoolbox 宏包动态修补 proof 环境，使其标题后自动换行
\patchcmd{\proof}{\item[\hskip\labelsep\itshape#1\@addpunct{.}]}{\item[\hskip\labelsep\itshape#1\@addpunct{.}]\hfill\par}{}{}

\begin{document}
\section*{\textbf{摘要}}

对于固定的 $h$，令 $N(h,k)$ 为满足以下条件的最小正整数 $N$：任意由 $k$ 个整数组成的集合所达到的基数 $|hA|$，亦必可由 $\{1,\dots,N\}$ 的某个 $k$ 元子集达到。借助 Rajagopal 关于 $|hA|$ 可能取值范围的固定 $h$ 定理，我们证明了对于任意固定的 $h$，均有 $N(h,k)\leq k^{C_h}$。本文的核心新贡献在于，利用多项式直径结构对 Rajagopal 大值构造中的稀疏几何块进行了替换。该替换操作在低次 Hilbert 函数层面实施，而低次 Hilbert 函数恰是 $h$ 固定时决定 $h$ 重和集性质的全部关键数据。
\section{引言}

对于有限集 $A \subset \mathbb{Z}$ 与正整数 $h$，记
\[
hA = \{a_1 + \cdots + a_h : a_i \in A\}
\]
其中允许元素重复。定义
\[
R(h,k) = \{|hA| : A \subset \mathbb{Z}, |A| = k\}.
\]
令 $N(h,k)$ 为满足以下条件的最小正整数 $N$：
\[
R(h,k) = \{|hA| : A \subseteq \{1,\dots,N\}, |A| = k\}.
\]
等价地，我们亦可考虑区间 $[0, N-1]$ 并在最后进行平移。

本文的目的是证明下述多项式压缩界。
\textbf{定理 1.1.} 对于任意固定的 $h \ge 1$，均存在常数 $C$，使得对一切 $k \ge 2$，成立
\[
N(h,k) \le k^{Ch}.
\]
本证明依赖于 Rajagopal 关于固定折叠和集大小取值范围的定理。令
\[
M = h^{k-h+1}, \quad K = \binom{h+k-1}{h,k},
\]
并定义
\[
\Delta = \bigcup_{\ell=0}^{\min(h,k)-3} [M_{h,k} + \ell h+1, M_{h,k} + \ell h + (h-2-\ell)].
\]
此处及后续内容中，$[u,v]$ 表示满足 $u \le n \le v$ 的整数 $n$ 构成的集合；当 $u > v$ 时，该区间视为空集。
\end{document}